\editionPage{sujet-A}{Sujet A · Automatismes}{Première partie · 6 points}{Corrigé p.~\pageref*{sujet-A-cor-auto}}
\begin{blocEdition}{sujet-A}
\textbf{Sujet original · Durée : 2 heures · L’usage de la calculatrice n’est pas autorisé.}
\par\smallskip\textbf{Première partie (6 points) : automatismes.} Pour chacune des 12 questions,
reporter son numéro sur la copie et donner une réponse courte, avec les unités utiles ;
aucune justification n’est demandée. Chaque question vaut 0,5 point.
\par\vspace{2mm}\begin{questions}\item Un prix de 60 euros augmente de $15\,\%$. Quel est le nouveau prix ?
\item Après une hausse de $25\,\%$, un article coûte 75 euros. Quel était son prix avant la hausse ?
\item Un prix baisse de $20\,\%$. Quel pourcentage de hausse permet de revenir au prix initial ?
\item Calculer $(2-\sqrt3)(2+\sqrt3)$.
\item Résoudre dans $\R$ l’équation $(2x-1)(x+3)=0$.
\item Résoudre dans $\R$ l’inéquation $-3x+6>0$.
\item Pour quels réels $x$ a-t-on $(x-1)(x+4)<0$ ?
\item Déterminer le coefficient directeur de la droite passant par $A(-1;3)$ et $B(3;-5)$.
\item On donne $P(A)=0{,}4$ et $P_A(B)=0{,}25$. Calculer $P(A\cap B)$.
\item Donner la médiane de la série $1\,;\,2\,;\,2\,;\,5\,;\,10$.
\item La figure 1 représente une fonction $f$ définie sur $[-2;4]$. Résoudre graphiquement $f(x)\leq1$.
\item Lire graphiquement l’équation réduite de la droite $d$ de la figure 2.\end{questions}
\label{figedition-sujet-A-1}\figureGeo{fr-sa-auto}
\end{blocEdition}

\editionPage{sujet-A-ex-1}{Sujet A · Exercice 1}{Deuxième partie · Une courbe et sa tangente}{7 points, dont 1 point de maîtrise de la langue\hfill Corrigé p.~\pageref*{sujet-A-cor-1}}
\begin{blocEdition}{sujet-A-ex-1}
\textbf{Deuxième partie (14 points) : deux exercices indépendants, notés chacun sur 7 points.}
\par\vspace{2mm}
On définit $f(x)=x^3-6x^2+9x+2$ pour $x\in[0;4]$.
\begin{questions}
\item \mbox{\textit{(1 point)}}\enspace Calculer $f'(x)$ et le factoriser.
\item \mbox{\textit{(2 points)}}\enspace Dresser le tableau de variations, avec les valeurs aux bornes et aux points critiques.
\item \mbox{\textit{(1 point)}}\enspace Donner le minimum et le maximum sur le segment, ainsi que toutes les abscisses où ils sont atteints.
\item \mbox{\textit{(1 point)}}\enspace Écrire la tangente $T$ au point d'abscisse 2.
\item \mbox{\textit{(1 point)}}\enspace Vérifier que $f(x)-(-3x+10)=(x-2)^3$, puis déterminer la position de la courbe par rapport à $T$.
\end{questions}
On attend un tableau lisible et une phrase de conclusion pour les extremums et la position relative.

\end{blocEdition}

\editionPage{sujet-A-ex-2}{Sujet A · Exercice 2}{Deuxième partie · Ateliers, défauts et gain}{7 points, dont 1 point de maîtrise de la langue\hfill Corrigé p.~\pageref*{sujet-A-cor-2}}
\begin{blocEdition}{sujet-A-ex-2}

L'atelier A fournit $60\,\%$ des articles et l'atelier B les autres.
Les taux de défaut sont respectivement $5\,\%$ et $10\,\%$.
On choisit un article au hasard. $D$ désigne l'évènement « l'article présente un défaut ».
\begin{questions}
\item \mbox{\textit{(1 point)}}\enspace Compléter l'arbre et calculer $P(D)$.
\item \mbox{\textit{(1 point)}}\enspace Sachant que l'article présente un défaut, calculer la probabilité qu'il provienne de B.
\item \mbox{\textit{(1 point)}}\enspace Le gain net est 12 unités sans défaut et $-8$ avec défaut. Donner sa loi.
\item \mbox{\textit{(1 point)}}\enspace Calculer l'espérance du gain.
\item \mbox{\textit{(1 point)}}\enspace Calculer sa variance.
\item \mbox{\textit{(1 point)}}\enspace Pour deux articles de défauts indépendants et de même loi, calculer la probabilité d'au moins un défaut.
\end{questions}

\label{figedition-sujet-A-ex-2-1}\figureGeo{fr-sa-arbre}
\end{blocEdition}

\editionPage{sujet-B}{Sujet B · Automatismes}{Première partie · 6 points}{Corrigé p.~\pageref*{sujet-B-cor-auto}}
\begin{blocEdition}{sujet-B}
\textbf{Sujet original · Durée : 2 heures · L’usage de la calculatrice n’est pas autorisé.}
\par\smallskip\textbf{Première partie (6 points) : automatismes.} Pour chacune des 12 questions,
reporter son numéro sur la copie et donner une réponse courte, avec les unités utiles ;
aucune justification n’est demandée. Chaque question vaut 0,5 point.
\par\vspace{2mm}\begin{questions}\item Calculer $\dfrac34+\dfrac16$.
\item Calculer $\dfrac{(2^3)^2}{2^4}$.
\item Factoriser $9x^2-1$.
\item Une quantité passe de 80 à 92. Calculer son taux d’évolution, en pourcentage.
\item Deux hausses successives d’un même taux donnent une hausse globale de $21\,\%$. Quel est ce taux ?
\item Exprimer $h$ en fonction de $V$ et de $r$, sachant que $V=\pi r^2h$ avec $r>0$.
\item La droite $d$ passe par $A(0;3)$ et a pour coefficient directeur $-2$. Donner son équation réduite.
\item Calculer la moyenne des notes 8, 12 et 15, de coefficients respectifs 1, 2 et 1.
\item \begin{minipage}[t]{86mm}Le tableau ci-contre répartit les 200 élèves d’un lycée. On choisit un élève au hasard ; $F$ : « être une fille » et $I$ : « être interne ». Calculer $P_F(I)$.\end{minipage}\hfill
\begin{minipage}[t]{62mm}\vspace{-3.2mm}\small\begin{tabular}[t]{@{}lccc@{}}\toprule
 & Interne & Externe & Total\\\midrule
Fille & 30 & 90 & 120\\
Garçon & 30 & 50 & 80\\
Total & 60 & 140 & 200\\\bottomrule
\end{tabular}\end{minipage}
\item Avec le même tableau, calculer $P(F\cap I)$, puis $P_I(F)$.
\item La figure 1 représente une fonction $g$ définie sur $[-2;4]$. Donner le signe de $g(x)$ selon $x$.
\item La figure 2 résume une série de notes. Donner sa médiane et son écart interquartile.\end{questions}
\label{figedition-sujet-B-1}\figureGeo{fr-sb-auto}
\end{blocEdition}

\editionPage{sujet-B-ex-1}{Sujet B · Exercice 1}{Deuxième partie · Une évolution avec apport}{7 points, dont 1 point de maîtrise de la langue\hfill Corrigé p.~\pageref*{sujet-B-cor-1}}
\begin{blocEdition}{sujet-B-ex-1}
\textbf{Deuxième partie (14 points) : deux exercices indépendants, notés chacun sur 7 points.}
\par\vspace{2mm}
On pose $u_0=80$ et $u_{n+1}=\frac34u_n+10$, pour $n\geq0$.
\begin{questions}
\item \mbox{\textit{(1,5 point)}}\enspace On pose $v_n=u_n-40$. Montrer que $v$ est géométrique et préciser ses paramètres.
\item \mbox{\textit{(1 point)}}\enspace En déduire l'expression de $u_n$.
\item \mbox{\textit{(1 point)}}\enspace Justifier que $u_n>40$ pour tout $n$ et que la suite décroît strictement.
\item \mbox{\textit{(1,5 point)}}\enspace On donne $3^4=81$, $4^4=256$, $3^5=243$ et $4^5=1024$.
Trouver le premier rang tel que $u_n\leq50$.
\item \mbox{\textit{(1 point)}}\enspace Écrire une boucle Python qui renvoie ce premier rang.
\end{questions}
La recherche du premier rang doit comporter une justification, pas seulement un essai réussi.

\end{blocEdition}

\editionPage{sujet-B-ex-2}{Sujet B · Exercice 2}{Deuxième partie · Cercle, tangente et projection}{7 points, dont 1 point de maîtrise de la langue\hfill Corrigé p.~\pageref*{sujet-B-cor-2}}
\begin{blocEdition}{sujet-B-ex-2}

Dans un repère orthonormé, on donne $A(0;0)$, $B(4;0)$, $C(1;3)$ et $M(5;4)$.
\begin{questions}
\item \mbox{\textit{(1,5 point)}}\enspace Déterminer les médiatrices de $[AB]$ et $[AC]$, puis leur point d'intersection $\Omega$.
\item \mbox{\textit{(1 point)}}\enspace En déduire l'équation du cercle circonscrit au triangle $ABC$.
\item \mbox{\textit{(0,5 point)}}\enspace Vérifier que $C$ appartient à ce cercle.
\item \mbox{\textit{(1,5 point)}}\enspace Déterminer la tangente au cercle en $C$.
\item \mbox{\textit{(1,5 point)}}\enspace Projeter orthogonalement $M$ sur cette tangente et calculer la distance correspondante.
\end{questions}

\label{figedition-sujet-B-ex-2-1}\figureGeo{fr-sb-geometrie}
\end{blocEdition}

